\documentclass[french,11pt]{book}
\usepackage{teach}
\usepackage[french]{babel}
\usepackage{amsmath,amssymb}
\usepackage{array}
\newcommand{\R}{\mathbb{R}}
\begin{document}
\begin{center}
{\Large \textbf{Devoir surveillé -- Suites numériques}}\\[2mm]
Terminale technologique
\end{center}

\section*{Devoir surveillé}
\begin{enumerate}
  \item Une entreprise vend 1200 objets en janvier. Les ventes augmentent de 4\% chaque mois. On note $u_n$ le nombre d'objets vendus $n$ mois après janvier. Donner $u_0$, la raison de la suite et calculer $u_6$.
  \item Un capital de 1500 euros est placé avec des intérêts composés au taux annuel de 2,5\%. Écrire la suite qui modélise le capital, puis calculer le capital au bout de 5 ans.
  \item Un club compte 240 adhérents. Chaque année, il perd 20 adhérents mais en gagne 35 nouveaux. Modéliser l'évolution par une suite et prévoir l'effectif après 8 ans.
\end{enumerate}

\end{document}
